\section*{الفصل \ref{ch:b1:counting} --- العدّ}

\begin{solution}{exo:b1:counting:1}
مراحل مستقلة وقاعدة الجداء: $26^2 \times 10^3 \times 26^2
= 26^4 \times 1000 = 456\,976\,000$ لوحة. وإذا كانت الحروف الأربعة
متمايزة مثنى مثنى كوّنت مراحل الحروف ترتيبةً من الرتبة $4$ في
الأبجدية: $26 \times 25 \times 24 \times 23 = 358\,800$ طريقة، ومنه
$358\,800 \times 1000 = 358\,800\,000$ لوحة.
\end{solution}

\begin{solution}{exo:b1:counting:2}
لكلمة \textsc{برتقال} $6$ حروف متمايزة: أي $6! = 720$ قلبًا.
ولكلمة \textsc{بانانا} $6$ حروف بتكرار ($3$ من الحرف ا، و $2$ من
الحرف ن، و $1$ من الحرف ب): فكل قلب محدَّد بمواضع حروف ا
($\binom 63$ اختيارًا)، ثم بمواضع حروف ن من بين المواضع $3$ الباقية
($\binom 32$)، ويأخذ الحرف ب الموضع الأخير:
$\binom{6}{3}\binom{3}{2} = 20 \times 3 = 60$ قلبًا
(أي، على نحو مكافئ، $6!/(3!\,2!\,1!) = 60$).
\end{solution}

\begin{solution}{exo:b1:counting:3}
الإجمالي: $\binom{12}{4} = 495$. وأمّا $2$ من النساء بالضبط: فنختارهما
($\binom 72 = 21$) ونختار $2$ من الرجال ($\binom 52 = 10$): أي $210$ لجنة.
وأمّا رجل واحد على الأقل: فهي متمّمة «لا رجل فيها»،
$\binom{12}{4} - \binom{7}{4} = 495 - 35 = 460$.
\end{solution}

\begin{solution}{exo:b1:counting:4}
\emph{شهور الميلاد:} الأدراج هي الشهور $12$؛ و $13$ شخصًا في $12$
درجًا يفرضون وجود شخصين في الدرج نفسه
(\cref{cor:b1:counting:pigeonhole}).

\emph{الأعداد المتتالية:} الأدراج هي الأزواج $n$ $\{1,2\},
\{3,4\}, \dots, \{2n-1, 2n\}$، وهي تجزّئ $\intint{1}{2n}$.
واختيار $n + 1$ عددًا صحيحًا يضع عددين في الزوج نفسه، وعنصرا
الزوج الواحد متتاليان.
\end{solution}

\begin{solution}{exo:b1:counting:5}
من أجل $1 \leq k \leq n$،
\[
k \binom nk = k\,\frac{n!}{k!\,(n-k)!}
= n\,\frac{(n-1)!}{(k-1)!\,(n-k)!} = n \binom{n-1}{k-1}.
\]
وبالجمع وإعادة الفهرسة بوضع $j = k - 1$:
\[
\sum_{k=0}^{n} k \binom nk = n \sum_{j=0}^{n-1} \binom{n-1}{j}
= n\, 2^{n-1}
\]
حسب \cref{prop:b1:counting:identities}. (وبطريقة أخرى: اشتقّ
$(1+x)^n = \sum_k \binom nk x^k$ وضع $x = 1$.)
\end{solution}

\begin{solution}{exo:b1:counting:6}
\omterm{def:b1:logic:map}{التطبيق} المتزايد تمامًا $f \colon \intint{1}{k} \to \intint{1}{n}$
محدَّد بصورته، وهي جزء من $\intint{1}{n}$ فيه $k$ عنصرًا (اسرد
الجزء بترتيب تزايدي)؛ وبالعكس فإن كل جزء من $k$ عنصرًا يعطي
تطبيقًا واحدًا بالضبط من هذا النوع. ومنه $\binom nk$ تطبيقًا متزايدًا تمامًا.

وإذا كان $f$ متزايدًا لا غير، فضع $g(i) = f(i) + i - 1$. عندئذ يكون $g$
متزايدًا تمامًا (فبين متغيرين متتاليين يربح $f$ مقدارًا $\geq 0$
ويربح $i - 1$ مقدار $1$) وقيمه في $\intint{1}{n + k - 1}$؛ و
$f(i) = g(i) - i + 1$ يستعيد $f$ من أيّ $g$ متزايد تمامًا
في $\intint{1}{n+k-1}$. وهذا تقابل، فيوجد إذن
$\binom{n + k - 1}{k}$ تطبيقًا متزايدًا.
\end{solution}

\begin{solution}{exo:b1:counting:7}
اقسم \omterm{def:b1:logic:sets}{مجموعة} $E$ فيها $m + n$ عنصرًا إلى كتلتين $M$ (وفيها $m$ عنصرًا)
و $N$ (وفيها $n$ عنصرًا). كل جزء من $E$ فيه $k$ عنصرًا يحوي $j$ عنصرًا
من $M$ (حيث $0 \leq j \leq k$) و $k - j$ عنصرًا من $N$؛ ومن أجل $j$ مثبَّت يوجد
$\binom mj \binom{n}{k-j}$ جزءًا كهذا، وتجزّئ الحالات $j = 0, \dots,
k$ الأجزاء ذات $k$ عنصرًا. وتعطي قاعدة الجمع متطابقة
فاندرموند\index{متطابقة فاندرموند}.

ومع $m = n = k$: $\binom{2n}{n} = \sum_{j=0}^{n} \binom nj
\binom{n}{n-j} = \sum_{j=0}^{n} \binom nj^2$، باستعمال $\binom{n}{n-j} =
\binom nj$.
\end{solution}

\begin{solution}{exo:b1:counting:8}
لتكن $A_d$ \omterm{def:b1:logic:sets}{مجموعة} مضاعفات $d$ في $\intint{1}{1000}$، ومنه
$\abs{A_d} = \lfloor 1000/d \rfloor$. وبالاحتواء والاستبعاد
(\cref{thm:b1:counting:inclexcl}) مع $A_2, A_3, A_5$، مع ملاحظة أن
$A_2 \cap A_3 = A_6$ وهكذا:
\[
500 + 333 + 200 - 166 - 100 - 66 + 33 = 734 .
\]
إذن يوجد $734$ عددًا صحيحًا يقبل القسمة على $2$ أو $3$ أو $5$.
\end{solution}

\begin{solution}{exo:b1:counting:9}
على \omterm{def:b1:logic:sets}{مجموعة} ذات $2$ عنصرين: جميع التطبيقات $2^4 = 16$ عدا التطبيقين
الثابتين $2$: أي $14$ تطبيقًا \omterm{def:b1:logic:inj}{شاملًا}.

وعلى \omterm{def:b1:logic:sets}{مجموعة} ذات $3$ عناصر: بالاحتواء والاستبعاد على القيم الفائتة،
يكون عدد التطبيقات من \omterm{def:b1:logic:sets}{مجموعة} ذات $4$ عناصر إلى \omterm{def:b1:logic:sets}{مجموعة} ذات $3$ عناصر
التي تفوتها قيمة واحدة على الأقل هو $\binom 31 2^4 - \binom 32 1^4 = 48 - 3 = 45$؛ وعدد التطبيقات كلها $3^4 =
81$؛ فعدد الشاملة: $81 - 45 = 36$. (وللتحقق: \omterm{def:b1:logic:inj}{التطبيق الشامل} من $4$ عناصر
على $3$ عناصر يكرّر قيمة واحدة بالضبط: فاختر القيمة المكررة
($3$)، والزوج الذي يُرسل إليها ($\binom 42 = 6$)، وتقابلًا من أجل
الباقي ($2$): $3 \times 6 \times 2 = 36$.)
\end{solution}

\begin{solution}{exo:b1:counting:10}
حلّ المعادلة $x_1 + \dots + x_n = k$ في $\N^n$ يُرمَّز بصفّ من
$k$ نجمة و $n - 1$ عصًا: اكتب $x_1$ نجمة، ثم عصًا، ثم $x_2$ نجمة، ثم
عصًا، \dots، وانتهِ بعدد $x_n$ من النجوم. وهذا تقابل على
كلمات الطول $k + n - 1$ المستعملة $k$ نجمة و $n - 1$ عصًا، وهذه
الكلمات محدَّدة بمواضع النجوم:
$\binom{n + k - 1}{k}$. والاختيارات بتكرار تقابل
حلول المعادلة (حيث $x_i$ عدد نسخ الغرض $i$)، فيكون العدد
نفسه.
\end{solution}

\begin{solution}{exo:b1:counting:11}
مع $A_i = \{\sigma : \sigma(i) = i\}$، تصمد \omterm{def:b1:counting:objects}{التبديلة} من
$\bigcap_{i \in I} A_i$ عند كل $i \in I$ وتبدّل النقاط
$n - \abs I$ الأخرى بحرية: $\abs{\bigcap_{i \in I} A_i} =
(n - \abs I)!$. وبالاحتواء والاستبعاد:
\[
\Bigl|\bigcup_i A_i\Bigr|
= \sum_{k=1}^{n} (-1)^{k+1} \binom nk (n-k)!
= \sum_{k=1}^{n} (-1)^{k+1} \frac{n!}{k!} ,
\]
لأن عدد أجزاء $I$ ذات الحجم $k$ هو $\binom nk$. ومنه
\[
D_n = n! - \Bigl|\bigcup_i A_i\Bigr|
= n!\Bigl(1 - \sum_{k=1}^{n} \frac{(-1)^{k+1}}{k!}\Bigr)
= n! \sum_{k=0}^{n} \frac{(-1)^k}{k!} .
\]
وأمّا المتطابقة الثانية: فصنّف التبديلات $\sigma$ في
$\intint{1}{n}$ بحسب \omterm{def:b1:logic:sets}{مجموعة} نقاطها الصامدة $F(\sigma)$. ومن أجل جزء
$F$ مثبَّت فيه $k$ عنصرًا، تكون التبديلات التي تحقق $F(\sigma) = F$ هي
بالضبط اضطرابات المتمّمة: أي $D_{n-k}$ \omterm{def:b1:counting:objects}{تبديلة}. وبالجمع على
اختيارات $F$ التي عددها $\binom nk$ من أجل كل $k$:
$n! = \sum_{k=0}^{n} \binom nk D_{n-k}$.
\end{solution}

\begin{solution}{exo:b1:counting:12}
\emph{المتطابقة الأولى.} عُدَّ الأزواج $(A, a)$ حيث $A \subseteq E$
(و $\abs E = n$) و $a \in A$. فبحسب حجم $A$: يوجد $\sum_k \binom nk k$
زوجًا. وباختيار العنصر المميَّز أولًا: يوجد $n$ اختيارًا للعنصر $a$، ثم
أيّ جزء من العناصر $n - 1$ الباقية لإتمام $A$:
أي $n\,2^{n-1}$ زوجًا.

\emph{المتطابقة الثانية.} عُدَّ الثلاثيات $(A, a, b)$ حيث $a, b \in
A$ (ويجوز $a = b$). فبحسب الحجم: $\sum_k k^2 \binom nk$. ومباشرةً:
إمّا $a = b$ (أي $n\,2^{n-1}$ ثلاثية، وهو العدّ السابق) وإمّا $a \neq b$
(أي $n(n-1)$ اختيارًا مرتَّبًا، ثم أيّ جزء من العناصر $n - 2$
الأخرى: $n(n-1)\,2^{n-2}$). والمجموع
\[
n\,2^{n-1} + n(n-1)\,2^{n-2} = n\,2^{n-2}\,(2 + n - 1)
= n(n+1)\,2^{n-2} .
\]
\end{solution}

\begin{solution}{pb:b1:counting:1}
\textbf{1.} $D_1 = 0$ (\omterm{def:b1:counting:objects}{فالتبديلة} الوحيدة تصمد عند $1$)، و $D_2 = 1$
(وهي التبديل)، و $D_3 = 2$ (بالترميز السطري: $231$ و $312$). ومن أجل
$n = 4$، بالتجميع بحسب $\sigma(1)$: مع $\sigma(1) = 2$ تكون الاضطرابات
$2143$ و $2341$ و $2413$؛ ومع $\sigma(1) = 3$: $3142$ و $3412$
و $3421$؛ ومع $\sigma(1) = 4$: $4123$ و $4312$ و $4321$. ثلاثة في كل
\omterm{def:b1:logic:sets}{مجموعة}: أي $D_4 = 9$.

\textbf{2.} \omterm{def:b1:counting:objects}{التبديلة} ذات $k$ نقطة صامدة بالضبط محدَّدة
باختيار \omterm{def:b1:logic:sets}{مجموعة} نقاطها الصامدة $F$ ($\binom nk$ طريقة) مع
قصرها على المتمّمة، الذي يجب أن يكون \omterm{def:b1:counting:objects}{تبديلة} على $n - k$ نقطة
\emph{بلا} نقطة صامدة ($D_{n-k}$ طريقة). والاختياران مستقلان
والتقارن تقابليّ: $P_k(n) = \binom nk D_{n-k}$.

\textbf{3.} لدينا $P_0(4) = D_4 = 9$؛ و$P_1(4) = \binom41 D_3 = 4 \times 2
= 8$؛ و$P_2(4) = \binom42 D_2 = 6$؛ و$P_3(4) = \binom43 D_1 = 0$
(إذ تفرض ثلاث نقاط صامدة نقطةً رابعة)؛ و $P_4(4) = 1$. والمجموع: $9 + 8 + 6
+ 0 + 1 = 24 = 4!$. فألا يوافق شيء ($9$ حالة) يفوق أن توافق رسالة واحدة بالضبط
($8$ حالات) --- بفارق ضئيل.

\textbf{4.} عُدَّ الأزواج $(\sigma, i)$ التي تحقق $\sigma(i) = i$. ومن أجل
$i$ مثبَّت، تكون التبديلات الصامدة عند $i$ هي تبديلات النقاط
$n - 1$ الأخرى: أي $(n-1)!$ \omterm{def:b1:counting:objects}{تبديلة}. ومنه فإن عدد الأزواج
هو $n \cdot (n-1)! = n!$، وهذا العدد يساوي كذلك
$\sum_\sigma \abs{\mathrm{Fix}(\sigma)}$. وبالقسمة على عدد
التبديلات $n!$: يكون متوسط عدد النقاط الصامدة مساويًا
$1$ بالضبط، من أجل كل $n \geq 1$.

\textbf{5.} ليكن $\sigma$ اضطرابًا في $\intint1{n+1}$ وليكن
$j = \sigma(n+1) \in \intint1n$: أي $n$ قيمة ممكنة. \emph{الحالة
$\sigma(j) = n+1$:} تتبادل النقطتان $j$ و $n+1$، ويكون قصر $\sigma$
على النقاط $n - 1$ الباقية اضطرابًا كيفيًا فيها:
أي $D_{n-1}$ إمكانًا. \emph{الحالة $\sigma(j)
\neq n+1$:} لتكن $i_0 = \sigma^{-1}(n+1)$؛ عندئذ $i_0 \neq j$ و
$i_0 \leq n$. نعرّف $\tau$ على $\intint1n$ بالعلاقتين $\tau(i) = \sigma(i)$
من أجل $i \neq i_0$ و $\tau(i_0) = j$. عندئذ يكون $\tau$ \omterm{def:b1:counting:objects}{تبديلة}
في $\intint1n$ (إذ حلّت القيمة الفائتة $j$ محلّ القيمة $n+1$)،
وهي اضطراب: إذ $\tau(i_0) = j \neq i_0$، و
$\tau(i) = \sigma(i) \neq i$ في المواضع الأخرى. وبالعكس، من
اضطراب $\tau$ في $\intint1n$ ومن القيمة $j$ نستعيد
$\sigma$ بوضع $\sigma(n+1) = j$ و$\sigma(\tau^{-1}(j)) = n+1$
و $\sigma = \tau$ في المواضع الأخرى: وهو تقابل يعطي $D_n$
إمكانًا. وبالجمع على $j$: $D_{n+1} = n(D_n + D_{n-1})$.
وعدديًا: $D_5 = 4(9 + 2) = 44$ و$D_6 = 5(44 + 9) = 265$.

\textbf{6.} من السؤال 5 لدينا $D_{n+1} = nD_n + nD_{n-1}$، ومنه
\[
u_{n+1} = D_{n+1} - (n+1)D_n = nD_n + nD_{n-1} - (n+1)D_n
= -(D_n - nD_{n-1}) = -u_n .
\]
وبما أن $u_1 = D_1 - 1 \cdot D_0 = -1$ فإن الاستقراء يعطي $u_n =
(-1)^n$، أي $D_n = nD_{n-1} + (-1)^n$ من أجل $n \geq 1$.

\textbf{7.} بالاستقراء على $n$. البداية: $D_0 = 1 = 0!\,s_0$. خطوة
الانتقال: بافتراض $D_{n-1} = (n-1)!\,s_{n-1}$،
\[
D_n = nD_{n-1} + (-1)^n = n!\,s_{n-1} + (-1)^n
= n!\Bigl(s_{n-1} + \frac{(-1)^n}{n!}\Bigr) = n!\,s_n ,
\]
وهي الصيغة المطلوبة. ولم يُستعمل أيّ احتواء واستبعاد: بل التراجع
التوفيقي من السؤال 5 وحده.

\textbf{8.} المراجعة الثلاثية، بالعوامل:
\[
\binom nk \binom kj
= \frac{n!}{k!\,(n-k)!} \cdot \frac{k!}{j!\,(k-j)!}
= \frac{n!}{j!\,(n-j)!} \cdot \frac{(n-j)!}{(k-j)!\,(n-k)!}
= \binom nj \binom{n-j}{k-j} .
\]
والآن عوّض $a_k = \sum_j \binom kj b_j$ وبادل بين المجموعين
المنتهيين:
\[
\sum_{k=0}^{n} (-1)^{n-k} \binom nk a_k
= \sum_{j=0}^{n} b_j \binom nj
\sum_{k=j}^{n} (-1)^{n-k} \binom{n-j}{k-j}
= \sum_{j=0}^{n} b_j \binom nj
\sum_{i=0}^{n-j} (-1)^{(n-j)-i} \binom{n-j}{i} .
\]
والمجموع الداخلي هو نشر $(1 + (-1))^{n-j} = 0^{n-j}$
(بمبرهنة ثنائي الحدّ، \cref{thm:b1:counting:binomial}): فهو ينعدم
من أجل $j < n$ ويساوي $1$ من أجل $j = n$. فلا يبقى إلا $j = n$، ويكون
الطرف الأيمن هو $b_n$، وهو المطلوب.

\textbf{9.} بالتناظر $\binom nk = \binom n{n-k}$، تُكتب
متطابقة \cref{exo:b1:counting:11} على الصورة $n! =
\sum_{k=0}^n \binom nk D_k$. وبتطبيق السؤال 8 مع $a_n = n!$ و
$b_k = D_k$:
\[
D_n = \sum_{k=0}^{n} (-1)^{n-k} \binom nk k!
= \sum_{k=0}^{n} (-1)^{n-k} \frac{n!}{(n-k)!}
= n! \sum_{j=0}^{n} \frac{(-1)^j}{j!} ,
\]
وبإعادة الفهرسة بوضع $j = n - k$: نحصل على الصيغة مرة ثالثة.

\textbf{10.} لدينا $D_n = n!\,s_n$ (السؤال 7)، ومنه
\[
\Bigl| D_n - \frac{n!}{\eu} \Bigr|
= n!\,\abs{s_n - \eu^{-1}} < \frac{n!}{(n+1)!} = \frac1{n+1} .
\]

\textbf{11.} من أجل $n \geq 1$ لدينا $\frac1{n+1} \leq \frac12$، وتكون
متراجحة السؤال 10 تامة: أي أن $D_n$ يقع على مسافة
$< \frac12$ من $n!/\eu$، فهو إذن أقرب عدد صحيح إليه على نحو وحيد.
ومن أجل $n = 0$ لا يعطي الحاصر إلا مسافة $< 1$، وتخفق الدعوى هناك
بالفعل: إذ إن $0!/\eu \approx 0.368$ أقرب عدد صحيح إليه $0$، بينما
$D_0 = 1$.

\textbf{12.} المقدار $\eu^{-1} - s_n = \sum_{k \geq n+1} (-1)^k/k!$ متسلسلة
متناوبة حدودها متناقصة تمامًا، فإشارته إشارة حدّه الأول
$(-1)^{n+1}/(n+1)!$. ومنه فإن إشارة $s_n -
\eu^{-1}$ هي إشارة $(-1)^n$: فمن أجل $n$ زوجيّ يكون $s_n > \eu^{-1}$
و$D_n = n!\,s_n > n!/\eu$؛ ومن أجل $n$ فرديّ يكون $D_n < n!/\eu$.

\textbf{13.} $D_7 = 6(265 + 44) = 6 \times 309 = 1854$؛ و$D_8 =
7(1854 + 265) = 7 \times 2119 = 14\,833$؛ و$D_9 = 8(14\,833 + 1854)
= 8 \times 16\,687 = 133\,496$؛ و$D_{10} = 9(133\,496 + 14\,833) =
9 \times 148\,329 = 1\,334\,961$. وللتحقق: $10!/\eu = 3\,628\,800 /
2.718281828 \approx 1\,334\,960.92$، وأقرب عدد صحيح إليه هو
$1\,334\,961$ --- ولدينا $D_{10} > 10!/\eu$، كما يتنبأ به السؤال 12
من أجل $n$ زوجيّ.

\textbf{14.} $\abs{p_n - \eu^{-1}} = \abs{s_n - \eu^{-1}} <
\frac1{(n+1)!}$. ومن أجل $n = 6$: $p_6 = 265/720 = 0.36806$ (بخمسة
أرقام عشرية)، في مقابل $\eu^{-1} = 0.36788$؛ والفارق دون $1/7! =
1/5040 < 2 \times 10^{-4}$. ويتقلص الحاصر $1/(n+1)!$ بسرعة تجعل
الاحتمال مثبَّتًا إلى أرقام عشرية كثيرة منذ اثنتي عشرة رسالة:
فالجواب «نحو $36.8\%$» مستقل عن $n$ في كل غرض عملي --- وهي مفاجأة
المسألة الشهيرة.

\textbf{15.} حسب السؤال 2 ومع $D_m = m!\,s_m$:
\[
\frac{P_k(n)}{n!} = \frac{\binom nk D_{n-k}}{n!}
= \frac{D_{n-k}}{k!\,(n-k)!} = \frac{s_{n-k}}{k!}
\;\longrightarrow\; \frac{\eu^{-1}}{k!}
\]
عندما $n \to \infty$ مع $k$ مثبَّت، لأن $s_{n-k} \to \eu^{-1}$.
والقيم الحدّية $\eu^{-1}/k!$ (حيث $k \in \N$) هي أوزان
توزيع بواسون ذي الوسيط $1$.

\textbf{16.} عُدَّ الثلاثيات $(\sigma, i, j)$ حيث $i \neq j$ و
$\sigma(i) = i$ و $\sigma(j) = j$. فباختيار الزوج المرتَّب أولًا:
$n(n-1)$ طريقة؛ والتبديلات الصامدة عند $i$ و $j$ معًا هي
تبديلات النقاط $n - 2$ الباقية: أي $(n-2)!$ \omterm{def:b1:counting:objects}{تبديلة}.
والمجموع: $n(n-1)(n-2)! = n!$. وأمّا الجمع على التبديلات $\sigma$ أولًا فيعدّ،
من أجل كل $\sigma$، الأزواج المرتَّبة من النقاط الصامدة
المتمايزة: $\abs{\mathrm{Fix}(\sigma)}(\abs{\mathrm{Fix}(\sigma)}-1)$.
ومنه المتطابقة المذكورة؛ وبالقسمة على $n!$ يكون متوسط
$\abs{\mathrm{Fix}}(\abs{\mathrm{Fix}} - 1)$ مساويًا $1$، فيكون متوسط
$\abs{\mathrm{Fix}}^2$ مساويًا $1 + 1 = 2$ ويكون التباين $2 -
1^2 = 1$.

\textbf{17.} النسب $1 - p_n$: من أجل $n = 4$، $1 - \frac
9{24} = \frac{15}{24} = 0.6250$؛ ومن أجل $n = 5$، $1 - \frac{44}{120} =
\frac{76}{120} = 0.6333$؛ ومن أجل $n = 6$، $1 - \frac{265}{720} =
\frac{455}{720} = 0.6319$. وكلها في حدود واحد في المئة من $1 - \eu^{-1}
\approx 0.6321$، متذبذبةً حوله.

\textbf{18.} مباشرةً:
\[
s_{n+2} - s_n = \frac{(-1)^{n+1}}{(n+1)!} +
\frac{(-1)^{n+2}}{(n+2)!}
= (-1)^{n+1}\Bigl(\frac1{(n+1)!} - \frac1{(n+2)!}\Bigr),
\]
والقوس $> 0$. فمن أجل $n$ زوجيّ يكون الفرق
سالبًا: أي $s_{n+2} < s_n$، ومنه $p_0 > p_2 > p_4 > \dots$؛ ومن أجل $n$
فرديّ يكون موجبًا: $p_1 < p_3 < p_5 < \dots$ وبالجمع مع
السؤال 12 (فالزوجية فوق $\eu^{-1}$ والفردية تحته) والسؤال 14
(فالمسافة إلى $\eu^{-1}$ تؤول إلى $0$): يحصر السلّمان
$\eu^{-1}$ بينهما.

\textbf{19.} السحبة الكاملة الواحدة \omterm{def:b1:counting:objects}{تبديلة} عشوائية منتظمة،
وتصحّ عندما تكون اضطرابًا: أي باحتمال $p_n \approx \eu^{-1}$.
وحسب الواقعة المذكورة، يكون متوسط عدد السحوب حتى النجاح
$1/p_n$، ويعطي السؤال 14 أن $1/p_n \approx \eu$ بخطأ
مهمَل منذ القيم الصغيرة للعدد $n$. إذن يكلّف تبادل الهدايا السرّي
مع إعادات السحب نحو $\eu \approx 2.72$ سحبة كاملة في المتوسط
--- سواء أكان في المكتب $6$ أشخاص أم $600$.

\textbf{20.} ثبّت $j \geq 2$ وأجرِ تصنيف السؤال 5 على
القيمة $\sigma(1) = j$. إذا كان $\sigma(j) = 1$: حملت النقاط
$n - 2$ الباقية اضطرابًا كيفيًا، أي $D_{n-2}$ طريقة. وإذا كان
$\sigma(j) \neq 1$: أعد توجيه السابقة $i_0 = \sigma^{-1}(1)$ إلى
$j$ تمامًا كما في السؤال 5؛ وهذا تقابل مع
اضطرابات النقاط $n - 1$ في $\{2, \dots, n\}$: أي $D_{n-1}$
طريقة. والمجموع $D_{n-1} + D_{n-2}$، وهو نفسه من أجل كل $j$. وبالجمع
على قيم $j$ التي عددها $n - 1$: $D_n = (n-1)(D_{n-1} + D_{n-2})$،
وهذا يُظهر العامل $n - 1$: أي $(n-1) \mid D_n$.

\textbf{21.} الدعوى: $D_n$ فرديّ إذا وفقط إذا كان $n$ زوجيًا. بالاستقراء
باستعمال $D_n = nD_{n-1} + (-1)^n$، أي $D_n \equiv nD_{n-1} + 1 \pmod 2$.
البداية: $D_1 = 0$ زوجيّ و $n = 1$ فرديّ: فتتحقق الدعوى. وإذا كان $n$ زوجيًا
كان $nD_{n-1}$ زوجيًا و $D_n \equiv 1$: أي فرديّ، وهو المطلوب. وإذا كان $n$
فرديًا كان $n - 1$ زوجيًا، فيكون $D_{n-1}$ فرديًا بالفرض، ويكون
$D_n \equiv D_{n-1} + 1 \equiv 0$: أي زوجيّ. وبهذا يُغلق الاستقراء.

\textbf{22.} إرجاع $D_n = nD_{n-1} + (-1)^n$ بترديد $n$ يُلغي
الحدّ الأول: $D_n \equiv (-1)^n \pmod n$. ومن أجل $n = 10$:
$(-1)^{10} = 1$، وينتهي $D_{10} = 1\,334\,961$ بالفعل بالرقم
$1$.

\textbf{23.} من أجل $n \geq 3$ لدينا $D_{n-1} \geq 1$، وقسمة
تراجع السؤال 6 على $D_{n-1}$ تعطي $D_n/D_{n-1} = n +
(-1)^n/D_{n-1}$، حيث $\abs{(-1)^n/D_{n-1}} \leq 1$ ويؤول
بسرعة إلى $0$. والاتساق: إذا كان $D_n \approx n!/\eu$ فإن
$D_n/D_{n-1} \approx n!/(n-1)! = n$ --- إذ يختصر العامل $\eu$
في النسبة، ويؤكد التراجع ذلك بدقة
$1/D_{n-1}$.

\textbf{24.} (أ) قاعدتا الجداء والجمع أساس كل عدّ:
فالسؤالان 2 و 5 يجزّئان مجموعات التبديلات إلى مراحل مستقلة.
(ب) وأعطى العدّ المزدوج المتوسط (السؤال 4) والتباين
(السؤال 16) لعدد النقاط الصامدة دون أيّ صيغة للعدد $D_n$ البتة.
(ج) وحسبت مبرهنة ثنائي الحدّ المجموع الداخلي المتناوب
$(1-1)^{n-j}$ الذي يجعل القلب الثنائي يعمل (السؤال 8). (د) وحوّل
حاصر المتسلسلة المتناوبة المجموعَ الدقيق لكن المعتم $n!\,s_n$ إلى
\omterm{def:b1:logic:statement}{العبارة} الشفافة «أقرب عدد صحيح إلى $n!/\eu$» (الأسئلة 10--14).

\textbf{25.} الاحتواء والاستبعاد
(\cref{ex:b1:counting:derangement} و\cref{exo:b1:counting:11})
هو البرهان المفهومي: فهو يشرح المجموع المتناوب بوصفه تصحيحات
لعدٍّ زائد، ويتعمّم حرفيًا على عدّ العناصر التي تتفادى أيّ عائلة
من المجموعات «السيئة». وأمّا طريق التراجع
(الأسئلة 5--7) فهو الأسرع حسابًا --- في زمن خطيّ، وبحساب صحيح
مضبوط، وبلا عوامل --- وهو مصدر الوقائع الحسابية في الجزء 5.
وأمّا القلب الثنائي (السؤالان 8--9) فيضع الصيغة داخل تحويل
عام سيعود كلما تقابل نسقان مثلثيان من المتطابقات. وأمّا
ظهور $\eu$ فأحسن ما يشرحه الصيغة نفسها: إذ إن نسبة
الاضطرابات هي المجموع الجزئي $s_n$ للمتسلسلة التي تعطي
$\eu^{-1}$، فكانت أظرفة مونمور، قبل ترميز أويلر بثلاثة عقود،
تحسب العدد $\eu$ فعلًا.
\end{solution}
